% evaluation.tex -- methodology section fragment.
% Intended to be \input into a paper that already loads amsmath (and, for the
% one external URL, hyperref or url). No preamble, no result tables.

\section{Evaluation}
\label{sec:evaluation}

We evaluate the concept-bottleneck model on four scenarios that probe
complementary capabilities: sparse point-annotation agreement on two
photo-quadrat benchmarks (CoralNet validation and Pacific Labeled Corals), dense
zero-shot benthic segmentation of orthomosaics (Benthos), and the linear
separability of the learned representation (CoralscapesV2 linear probe). All four
are driven by a single reproducible entry point,
\texttt{scripts/evaluate.py}, which loads one checkpoint and runs any subset of
the evaluations at the source (native) resolution of each dataset. This section
specifies the shared inference protocol, the ground-truth label mapping, the
metric definitions, and each scenario precisely enough to reproduce the
pipeline.

\subsection{Shared inference protocol}
\label{sec:eval-protocol}

\paragraph{Model.}
Unless stated otherwise the model is the concept-bottleneck segmenter
\texttt{ConceptBottleneckDPTLoRADINOv3}
(\texttt{configs/model\_config\_cbm\_dpt\_lora\_vitl.yaml}). A frozen DINOv3
ViT-L/16 backbone with LoRA adapters ($r=16$, $\alpha=32$, dropout $0.05$, applied
to the attention $q,k,v,o$ projections) feeds a DPT head that reassembles four
intermediate token blocks into a multi-scale pyramid and produces a
$256$-dimensional dense feature map. This feature map is bilinearly upsampled to
the model input resolution and passed through a $1{\times}1$ convolution (the
concept projection) that yields $K=704$ concept logits per pixel. The concept
logits are activated group-wise: for each of the six taxonomic ranks
(\texttt{kingdom}, \texttt{phylum}, \texttt{class}, \texttt{order},
\texttt{family}, \texttt{genus}; $6,16,28,74,176,367$ channels respectively,
$667$ channels in total) a softmax is taken over that rank's channel group, and a
sigmoid is applied independently to the remaining $37$ binary channels ($24$
morphology, $2$ health, $11$ non-coral).\footnote{%
\texttt{mermaidseg/model/models.py}, method \texttt{concept\_outputs\_activation}.}
The activated concept map (all values in $[0,1]$) is the model's bottleneck; a
final $1{\times}1$ convolution (the class head) maps it to $C=79$ MERMAID class
logits, where class id $0$ denotes \emph{ignore}/background ($78$ foreground
classes). The evaluation reads two quantities from a forward pass: the
per-pixel class probabilities $\mathrm{softmax}$ over the class logits, and the
activated concept map.\footnote{%
\texttt{mermaidseg/evaluation/predictor.py}, \texttt{CBMPredictor.forward\_probs}.}

\paragraph{Preprocessing.}
Each source image is resized (anisotropically) to the model input size
$512\times512$ and normalised with ImageNet statistics
($\mu=(0.485,0.456,0.406)$, $\sigma=(0.229,0.224,0.225)$); this reproduces the
validation transform used during training. No test-time augmentation, cropping,
or tiling of the network input is applied. The encoder optionally runs under
\texttt{bfloat16} autocast; all reported quantities (probabilities, features)
are cast to \texttt{float32} before scoring.\footnote{%
\texttt{make\_eval\_transform} and \texttt{CBMPredictor} in
\texttt{predictor.py}.}

\paragraph{Source-resolution scoring.}
All metrics are computed at each dataset's native resolution rather than at the
$512\times512$ network resolution, so that resizing never discards or invents
ground-truth structure. Concretely, a prediction map defined on the
$512\times512$ grid is bilinearly upsampled (\texttt{align\_corners=False}) to the
source resolution $H\times W$. For the point-annotation scenarios this upsample
is evaluated \emph{exactly} at the annotated pixels via \texttt{grid\_sample},
using the pixel-centre convention
\begin{equation}
x = \frac{c+0.5}{W}\cdot 2 - 1,
\qquad
y = \frac{r+0.5}{H}\cdot 2 - 1,
\label{eq:gridsample}
\end{equation}
for an annotation at source row $r$ and column $c$; this is identical to
bilinearly upsampling the full map to $H\times W$ and indexing $[\,r,c\,]$, but
avoids materialising the full-resolution map.\footnote{%
\texttt{CBMPredictor.sample\_at\_points} / \texttt{upsample\_scores}.} For the
dense scenarios the per-class score maps are upsampled to $H\times W$ and then
argmaxed.

\paragraph{Data access.}
CoralNet, Pacific Labeled Corals, and Benthos are read from an S3 mirror through
a local cache configured from the data config
(\texttt{setup\_local\_cache}); CoralscapesV2 is read either from the
HuggingFace dataset \texttt{josauder/coralscapesV2} or from a local
Cityscapes-style mirror.

\subsection{Mapping source labels to model outputs}
\label{sec:eval-mapping}

The benchmarks use heterogeneous label vocabularies, so both the class ground
truth and the concept ground truth are derived from each source label by
deterministic lookups that reproduce the training-time semantics.\footnote{%
\texttt{mermaidseg/evaluation/gt\_mapping.py}.}

\paragraph{Class ground truth.}
Each source label name is first mapped to a MERMAID benthic-attribute name via a
per-source map (a static JSON file
\texttt{configs/coralnet\_to\_mermaid\_mapping\_temporary.json} for CoralNet; a
static in-code dictionary for Pacific Labeled Corals). The attribute name is then
\emph{rolled up} through the MERMAID benthic-attribute parent hierarchy
(fetched from \texttt{https://api.datamermaid.org/v1/benthicattributes/} and
cached to \texttt{benthic\_hierarchy.json}) until it lands on a name present in
the model's frozen class vocabulary \texttt{id2label}; if no ancestor is in the
vocabulary the point maps to class $0$ and is ignored for class metrics. This
matches training with \texttt{label\_roll\_up=True}
(\texttt{registry.roll\_up\_label}).

\paragraph{Concept ground truth.}
Each source label is mapped to a length-$K$ concept row aligned to the model's
concept channel order, read from \texttt{configs/class\_to\_concepts.csv} for the
matching \texttt{(source\_dataset\_source, source\_label\_class\_name)} pair
(first row wins on duplicates; labels absent from the CSV map to an all-zero row
that is ignored everywhere). Each channel takes a value in
$\{0,1,2\}$ with the interpretation
\emph{$0$ = not given}, \emph{$1$ = inactive / False}, \emph{$2$ = active /
True}. For a taxonomic rank, a \texttt{not\_given}/empty cell sets the whole rank
to $0$; otherwise the channel matching the cell value is set to $2$ (active) and
the sibling channels of that rank to $1$ (inactive). For a binary concept, cells
\texttt{true}/\texttt{false} give $2$/$1$ and anything else gives $0$.\footnote{%
\texttt{build\_class\_lookup}, \texttt{build\_concept\_lookup},
\texttt{\_encode\_concept\_row} in \texttt{gt\_mapping.py}.}

\subsection{Metrics}
\label{sec:eval-metrics}

All accumulators are additive integer count tensors, so a \emph{pooled} (micro)
aggregate is the element-wise sum of the per-unit accumulators, while a
\emph{macro} aggregate is the unweighted mean of the per-unit scalar metrics.\footnote{%
\texttt{mermaidseg/evaluation/metrics.py}.}

\paragraph{Class accuracy and mIoU.}
From a confusion matrix $M$ (rows = ground truth, columns = prediction) built
over the evaluated points/pixels, with ground-truth class $0$ excluded and
predictions clipped into range,
\begin{equation}
\mathrm{acc} = \frac{\sum_c M_{cc}}{\sum_{c,c'} M_{cc'}},
\qquad
\mathrm{IoU}_c = \frac{M_{cc}}{\sum_{c'} M_{cc'} + \sum_{c'} M_{c'c} - M_{cc}} .
\end{equation}
The mean IoU is the average of $\mathrm{IoU}_c$ over the evaluated classes whose
union is non-zero, excluding the ignore class.\footnote{\texttt{ClassConfusion}.}

\paragraph{Taxonomic rank accuracy.}
For a rank $\rho$ with concept channel set $S_\rho$, a point is \emph{given} at
rank $\rho$ if some channel in $S_\rho$ has ground-truth value $2$; the
ground-truth class is that channel, and the prediction is
$\arg\max_{i\in S_\rho}$ of the model's activated concept probabilities. We
report \texttt{acc\_all} over all given points and \texttt{acc\_living} over the
subset whose ground-truth channel is not the rank's \texttt{$\rho$\_\_none}
channel.\footnote{\texttt{TaxonomicRankAccuracy}.}

\paragraph{Binary concept accuracy and F1.}
A binary concept is scored only where its ground truth is given (value in
$\{1,2\}$); the prediction is $p>0.5$ on the sigmoid-activated probability, and
\texttt{True} ($=2$) is the positive class. Per concept,
$\mathrm{acc}=(\mathrm{TP}+\mathrm{TN})/(\mathrm{TP}+\mathrm{FP}+\mathrm{FN}+\mathrm{TN})$
and $\mathrm{F1}=2\,\mathrm{TP}/(2\,\mathrm{TP}+\mathrm{FP}+\mathrm{FN})$. The
reported macro accuracy averages over concepts with at least one scored point,
and the macro F1 averages over concepts with $2\,\mathrm{TP}+\mathrm{FP}+\mathrm{FN}>0$.\footnote{\texttt{BinaryConceptStats}.}

\paragraph{Aggregation across units.}
The two point-annotation scenarios partition their images into \emph{units}
(defined per scenario below), compute the metrics above per unit, and report
both the pooled (micro, summed counts) and macro (mean of per-unit scalars)
aggregates.\footnote{\texttt{run\_point\_evaluation} in
\texttt{point\_datasets.py}.}

\subsection{Scenario 1: CoralNet validation}
\label{sec:eval-coralnet}

The validation split is a fixed whitelist of $21$ CoralNet sources
(\texttt{data\_config\_512.yaml}: source ids
\texttt{1645, 1776, 1783, 3058, 3466, 3573, 4614, 4721, 5155, 5181, 5444, 5847,
6537, 6699, 6700, 6721, 6834, 6838, 6918, 6926, 6968}), disjoint from the $42$
training sources; the entire annotated content of each whitelisted source is
used (no additional per-image sub-sampling). Only \emph{Confirmed} point
annotations are retained (filtered during data ingestion), and annotation
coordinates are expressed in the stored image space (longest edge clamped to
$2048$\,px), which the source-resolution sampling of
Eq.~\eqref{eq:gridsample} respects. The evaluation unit is the CoralNet
\texttt{source\_id}; we report per-source and pooled/macro class accuracy and
mIoU, per-rank taxonomic accuracy, and binary-concept macro accuracy/F1.\footnote{%
\texttt{mermaidseg/evaluation/coralnet\_eval.py}.}

\subsection{Scenario 2: Pacific Labeled Corals}
\label{sec:eval-pacific}

The Pacific Labeled Corals benchmark comprises three sites --- \texttt{heron\_reef},
\texttt{line\_islands}, and \texttt{nanwan\_bay} --- across the \texttt{reference} and
\texttt{evaluation} subsets. This dataset is never used for training
(\texttt{train: None}), so it is a fully held-out benchmark. Each point's source
label is taken from the primary annotator column \texttt{host}, falling back to
the \texttt{archived} annotation where \texttt{host} is missing, and points with
no label after fallback are dropped. Annotation \texttt{row}/\texttt{col} are the
native pixel coordinates supplied with the dataset. The evaluation unit is the
site, and the reported metrics are identical to
Scenario~1.\footnote{\texttt{mermaidseg/evaluation/pacific\_eval.py}.}

\subsection{Scenario 3: Benthos zero-shot benthic segmentation}
\label{sec:eval-benthos}

The Benthos benchmark evaluates dense, zero-shot benthic segmentation on two
orthomosaics: \texttt{RS24} (Red Sea, $11070\times11786$\,px) and
\texttt{CR\_DoubleWreck} (Caribbean, $24192\times3982$\,px). Each mosaic is cut
into non-overlapping $2048\times2048$ tiles in row-major order (edge tiles may be
smaller; no padding); each tile is processed independently through the shared
protocol of Section~\ref{sec:eval-protocol}.

Because the model does not output the Benthos label taxonomy directly, the
$8$ target classes are defined \emph{zero-shot} as per-pixel expressions over the
activated concept probabilities, via an editable specification
(\texttt{configs/eval/benthos\_zero\_shot.yaml}):
\begin{itemize}
\item Hard Coral $=$ \texttt{class\_\_hexacorallia * live};
\item Soft Coral $=$ \texttt{class\_\_octocorallia * live};
\item Sand $=$ \texttt{sand};
\item Rock $=$ \texttt{hard\_substrate};
\item Sponge $=$ \texttt{phylum\_\_porifera};
\item Other $=$ \texttt{anthropogenic};
\item Algae $=$ \texttt{algae * (1 - calcifying)};
\item Calcareous Algae $=$ \texttt{calcifying * algae}.
\end{itemize}
Expressions are parsed by a small per-pixel DSL in which \texttt{*}, \texttt{+},
and \texttt{-} multiply, add, and subtract concept probability maps with the
result clamped to $[0,1]$, numeric literals are constant maps, and
\texttt{rank:value} is shorthand for the \texttt{rank\_\_value} channel.\footnote{%
\texttt{mermaidseg/model/concept\_expr.py}; the driver is
\texttt{mermaidseg/evaluation/benthos\_eval.py}.} The specification is validated
against the checkpoint at run time: every concept atom it references (here
including \texttt{calcifying}, \texttt{algae}, \texttt{live},
\texttt{class\_\_hexacorallia}, \texttt{class\_\_octocorallia},
\texttt{phylum\_\_porifera}, \texttt{sand}, \texttt{hard\_substrate},
\texttt{anthropogenic}) must be present in the model's concept vocabulary, so the
spec must be matched to the checkpoint being evaluated.

For each tile, the eight expression maps are evaluated on the concept
probabilities, bilinearly upsampled to the tile's source resolution, and
argmaxed to a dense prediction. Nodata pixels (the transparent orthomosaic
border, RGB $=0$) and pixels whose ground-truth class is \texttt{unlabeled} or
not listed in the spec are excluded. We report pixel accuracy, per-class IoU,
and mIoU over the eight-way confusion matrix (which has no separate ignore
class), per site and pooled across sites.

\subsection{Scenario 4: CoralscapesV2 linear probe}
\label{sec:eval-coralscapes}

The final scenario measures the linear separability of the learned
representation on CoralscapesV2, a dense benthic-segmentation dataset with $95$
foreground classes (native label $0$ is background/ignore) at
$1024\times2048$\,px.\footnote{%
\texttt{mermaidseg/evaluation/coralscapes\_probe.py}; probe set built by
\texttt{scripts/build\_coralscapes\_v2\_probe\_points.py}.}

\paragraph{Fixed probe set.}
To make the probe reproducible and comparable across models, the training pixels
used to fit the probe are pinned in
\texttt{configs/eval/coralscapes\_v2\_probe\_points.json}. The builder
(\texttt{scripts/build\_coralscapes\_v2\_probe\_points.py}) targets $10$ points
per class and is deterministic (sorted paths, \texttt{ndimage.label}, and
\texttt{numpy.argmax} on the distance transform; there is no random seed). It
censuses $8$-connected components of at least $4$\,px on the full training
split. When a class has at least $10$ such components, $10$ are chosen by
round-robin across images ordered by their largest component, one interior pixel
each. When it has $N<10$, every component is used and the $10$ points are split
evenly, the largest components receiving the remainder (four components yield
quotas $3,3,2,2$). A quota of one is the distance-transform maximum. A larger
quota repeats that maximum and then removes every pixel with
$dy^{2}+dx^{2}\le 16^{2}$. If that disk empties a component before the quota is
met, the pixels already placed are kept and the unused quota is not moved onto
another component, so the class is stored with fewer than $10$ points. The
shortfall is recorded in the JSON. A class with no component of at least
$4$\,px, a component with fewer pixels than its quota, or a component that
yields no pixel still aborts the write.
The image list is exactly the images that own a probe pixel, each required to be
$1024\times2048$. Probe images are located either by file stem in a local mirror
or by ground-truth label MD5 in the HuggingFace dataset. The JSON header records
\texttt{num\_images}, \texttt{total\_points}, \texttt{points\_per\_class\_target},
and \texttt{placement\_shortfalls}.
On the local mirror of $1790$ training masks, three components can hold only one
probe pixel before the radius-$16$ disk removes the rest: \texttt{seriatopora dead}
components of $251$\,px and $156$\,px in \texttt{site30/site30\_000148\_000236}
(each quota $3$, one pixel kept) and a \texttt{turbinaria dead} component of
$65$\,px in \texttt{site33/site33\_000058\_033600} (quota $2$, one pixel kept).
Those two classes therefore have $6$ and $9$ probe pixels. Every other class has
$10$. The pinned set has $663$ images and $945$ points.

\paragraph{Strided prediction.}
Each $1024\times2048$ image is scored with three windows, stacked as one batch
and resized to the model input size ($256$ or $512$): the left half (columns
$0$:$1024$), the center square (columns $512$:$1536$), and the right half
(columns $1024$:$2048$). Raw concept logits are bilinearly upsampled back onto
those rectangles on the GPU and averaged with equal weight where the windows
overlap (columns $0$:$512$ see the left window only, $512$:$1024$ the left and
center, $1024$:$1536$ the center and right, $1536$:$2048$ the right only). The
model's concept activation --- softmax within each taxonomic rank, sigmoid on
the binary tail --- is applied to the averaged logits. An image that is not
$1024\times2048$ raises. Only the \texttt{concepts} feature kind is accepted;
\texttt{dpt} and \texttt{backbone} raise.

\paragraph{Probe fitting and dense evaluation.}
Probe pixels are read from that activated map at the stored source-resolution
coordinates, so the SVM sees the same features as the dense test. A linear probe
is fit with a standardiser followed by a linear SVM
(\texttt{StandardScaler} then \texttt{LinearSVC}, $C=1$,
\texttt{max\_iter}=$20000$). For dense inference the standardiser is folded into
the SVM's linear layer, giving effective weights
$W_\mathrm{eff}=W/\sigma$ and bias
$b_\mathrm{eff}=b - (\mu/\sigma)\,W^{\!\top}$, applied in row strips on the GPU
to the $1024\times2048$ activated concept map and argmaxed. We report the probe
train accuracy and, on the CoralscapesV2 test split, accuracy and mIoU over
classes $1$--$95$ (label $0$ ignored).

\paragraph{Bleached.}
Separately from the SVM, the sigmoid of the averaged \texttt{bleached} logit is
thresholded at $0.5$ (exactly $0.5$ counts as negative). Ground truth comes from
\texttt{configs/class\_to\_concepts.csv}: \texttt{TRUE} (code $2$) is the
positive class, \texttt{FALSE} (code $1$) is the negative class, and
\texttt{not\_given} (code $0$) is excluded. Mask id $0$ is unannotated and is
excluded as well; it is not the class named \texttt{background}, which is an
explicit \texttt{FALSE}. Alive, dead, and algae-covered corals are \texttt{FALSE}
when the table says so, so predicting bleached there is a false positive. A
native class missing from the table, or a cell other than \texttt{TRUE},
\texttt{FALSE}, or \texttt{not\_given}, raises. We report accuracy, precision,
recall, F1, and the TP/FP/FN/TN counts, plus how many pixels were
\texttt{not\_given} or unannotated. The same four scores are reported for
CoralNet points, where code $0$ (including the $23$ CoralNet labels whose
bleached cell is \texttt{not\_given}) is excluded rather than treated as
\texttt{FALSE}.

\subsection{Reproducing the evaluation}
\label{sec:eval-repro}

All four scenarios are produced by a single command; for example
\begin{quote}\ttfamily
uv run python scripts/evaluate.py {-}{-}checkpoint <ckpt> \\
\quad {-}{-}model-config configs/model\_config\_cbm\_dpt\_lora\_vitl.yaml \\
\quad {-}{-}data-config configs/data\_config\_512.yaml \\
\quad {-}{-}output-dir <out> {-}{-}evals all {-}{-}batch-size 8 {-}{-}num-workers 16 {-}{-}amp
\end{quote}
The required assets are the model checkpoint and its config, the class vocabulary
\texttt{demo/id2label.json}, the concept vocabulary
\texttt{demo/concept\_id2name.json} (which must match the checkpoint's concept
channels), the concept table \texttt{configs/class\_to\_concepts.csv}, the
Benthos spec \texttt{configs/eval/benthos\_zero\_shot.yaml}, the CoralscapesV2
probe set \texttt{configs/eval/coralscapes\_v2\_probe\_points.json}, and the data
config that selects the validation whitelists. The run writes a top-level
\texttt{summary.json}/\texttt{summary.md} and, per scenario, a
\texttt{metrics.json} (plus \texttt{per\_unit.csv} for the point scenarios) that
are flushed incrementally; the benthic hierarchy is cached to
\texttt{benthic\_hierarchy.json} on first use. Debug limiters
(\texttt{{-}{-}max-images-per-unit}, \texttt{{-}{-}max-tiles-per-site},
\texttt{{-}{-}max-test-images}) restrict the workload for smoke tests but are not
used for reported numbers. Given a fixed checkpoint and the pinned assets above,
the pipeline is deterministic: preprocessing, sampling, and the folded linear
probe involve no stochastic augmentation, and the SVM is fit on the fixed probe
point set.
